\documentclass[a4paper, 11pt]{article}

% -- Coments
\usepackage{verbatim}

% ----- Fonts -----
% -- Fuente --
% \usepackage{fontspec}
% \setmonofont{JetBrainsMono Nerd Font}

\usepackage{babel}
\babelprovide[import, main]{spanish}

% -- Color --
\usepackage{xcolor}
%\definecolor{azul}{RGB}{00,33,99}
\definecolor{azul}{RGB}{35,72,180}

% -- Page Margin --
\usepackage[margin=1in, headheight=15pt]{geometry}

% -- Encabezado con nombre y expediente --
\usepackage{fancyhdr}
\pagestyle{fancy}
\fancyhf{}
\lhead{\textbf{Bórquez Guerrero Angel Fernando}}
\rhead{Expediente: 219208106}
\renewcommand{\headrulewidth}{0.5pt}

% -- Espaciados --
\newcommand{\Pspace}{0.5cm}
\newcommand{\Aspace}{0.2cm}

% -- Columnas --
\usepackage{multicol}

% -- Imagenes --
\usepackage{graphicx}
\usepackage{float}

% -- Matemáticas --
\usepackage{amsmath, amssymb}

% -- Gráficas --
\usepackage{pgfplots}
\pgfplotsset{compat=1.18}


\title
{
    \textbf{Contexto Regional 2026-2} \\[0.3cm]
    {\large \textit{Importancia de los inicios de la maquila en la conformación industrial y laboral de Nogales, Sonora}} \\[0.2cm]
    {\normalsize Quintero Ramírez, Cirila} \\[0.3cm]
    {\normalsize Profesor José Moreno Vega}
}
\date{}

\begin{document}

\begin{otherlanguage}{spanish}
\maketitle
\end{otherlanguage}
\thispagestyle{fancy}
\vspace{-0.5cm}


{\Large 1. ¿Cuál es el argumento central del autor?} \\
{\color{azul}
    La autora sostiene que los inicios de la maquila en Nogales (1963-1975) dependieron sobre todo del apoyo del gobierno estatal. Desde el Plan de Diez Años de 1963, las leyes de fomento industrial y la creación del Parque Industrial Nogales, se ofrecieron exenciones de impuestos, terrenos y otras facilidades para atraer inversión extranjera, como la de Motorola. En un principio se pensó usar la maquila como base de un desarrollo industrial regional, pero esas mismas facilidades hicieron que la industria dependiera de los apoyos y de empresas extranjeras, sin vincularse con la economía local. Por eso, cuando cambiaron los ciclos de Estados Unidos, el mercado laboral se volvió inestable y llegó la primera crisis de 1974-1975, que mostró que la maquila creaba empleos pero también era frágil.
}
\vspace{1cm}

{\Large 2. ¿Qué evidencias utiliza para respaldarlo?} \par
{\color{azul}
Se apoya sobre todo en números: cuadros con datos del INEGI y de la Secretaría de Programación y Presupuesto que muestran cómo los empleos de la maquila pasaron de 18 en 1967 a 9,640 en 1974, y cómo estaba repartida la gente que trabajaba en Nogales por sector y por sexo. También revisa las leyes y los apoyos que dio el gobierno, como la Ley 16 de 1962, la Ley 49 de 1967 y el convenio de 1969 para el parque industrial, y pone casos concretos como el de Motorola y el programa \textit{shelter} de Richard Campbell. Además usa periódicos de la época, como \textit{El Sol} y \textit{La Lucha}, discursos del presidente Echeverría y trabajos de otros autores como Kopinak, Cravey y José Carlos Ramírez.
}
\vspace{1cm}

{\Large 3. ¿Cómo se relaciona esta lectura con la lectura anterior?} \par
{\color{azul}
Las dos hablan de Sonora en los años sesenta, cuando los gobiernos impulsaron proyectos de desarrollo con la idea de traer progreso, la presa El Novillo en un caso y la maquila en el otro. En ambas lecturas el gobierno estatal tiene un papel central, y la gente común queda con los costos: los desplazados fueron excluidos de las indemnizaciones, y los obreros quedaron expuestos a los despidos de la crisis de 1974.
}

\end{document}
